\chapter{Write the Black-Box Test}

The manual request proved that the test jail can observe the rendered product
through nginx. We will now make that contract executable with Ruby and
Minitest. The runner does not need Rails: Rails is part of the system under
test, while the test itself is an ordinary HTTP client.

\section{Exercise 22A: Install the test toolchain}

In the test jail, install FreeBSD's Ruby and Minitest packages:

\begin{lstlisting}[language=sh]
pkg install -y ruby rubygem-minitest
\end{lstlisting}

Verify both pieces without creating any files:

\begin{lstlisting}[language=sh]
ruby --version
ruby -rminitest -e 'puts Minitest::VERSION'
\end{lstlisting}

Exact patch versions may advance with the package repository. Both commands
must exit successfully; no Rails, database client, compiler, or browser driver
is needed in this jail.

\section{Exercise 22B: Express the external contract}

Create a directory and test file:

\begin{lstlisting}[language=sh]
mkdir -p /usr/local/tests
vi /usr/local/tests/store_e2e_test.rb
\end{lstlisting}

Enter:

\begin{lstlisting}[language=Ruby]
require "minitest/autorun"
require "net/http"
require "uri"

class StoreEndToEndTest < Minitest::Test
  URL = ENV.fetch("STORE_URL", "http://10.17.89.26/")

  def test_product_is_visible_through_the_web_edge
    uri = URI(URL)
    response = Net::HTTP.start(
      uri.host,
      uri.port,
      open_timeout: 2,
      read_timeout: 5
    ) do |http|
      http.get(uri.request_uri)
    end

    assert_equal "200", response.code
    assert_includes response.body, "FreeBSD Handbook"
    assert_includes response.body, "$49.99"
  end
end
\end{lstlisting}

The default endpoint is nginx, not store's Puma port. The assertions therefore
cross nginx, store, catalog, and PostgreSQL. \texttt{STORE\_URL} keeps the test
portable: another environment can point the same contract at a different web
edge without editing the test.

\section{Exercise 22C: Run the test}

Execute it directly with Ruby:

\begin{lstlisting}[language=sh]
ruby /usr/local/tests/store_e2e_test.rb
\end{lstlisting}

Expect one run, no failures, and no errors. Minitest may count more assertions
than the three method calls because \texttt{assert\_includes} checks both the
container and its content.

Run it once more with the endpoint made explicit:

\begin{lstlisting}[language=sh]
STORE_URL=http://10.17.89.26/ \
  ruby /usr/local/tests/store_e2e_test.rb
\end{lstlisting}

This second run should produce the same result and demonstrates the input a CI
job would supply.

\section{Checkpoint}

Stop after both invocations pass. We now have an automated black-box test over
the complete application stack. The next exercise will deliberately break one
dependency and confirm that the test fails for an observable reason rather
than passing accidentally.
